% ------------------------------------------------------------
% Information-Geometric Correspondence (IGC): Extended Handout
% ------------------------------------------------------------
\documentclass[12pt]{article}
\usepackage{amsmath,amssymb,amsthm}
\usepackage{geometry}
\usepackage{mdframed}
\usepackage{enumitem}
\usepackage{hyperref}
\usepackage[table]{xcolor}


\geometry{margin=1in}

% theorem environments
\newtheorem{definition}{Definition}
\newtheorem{theorem}{Theorem}
\newtheorem{corollary}{Corollary}
\newtheorem{remark}{Remark}
\pagecolor{white}

% boxed environment for emphasis
\newmdenv[linecolor=black,linewidth=0.5pt,backgroundcolor=gray!0,roundcorner=5pt]{mybox}

\begin{document}

\begin{center}
{\LARGE \textbf{Information-Geometric Correspondence (IGC)}}\\[0.5em]
{\large A Field Guide to Excess Fisher Curvature Across Domains}
\end{center}

\vspace{1em}

% --------------------------
\section*{1. Core Principle}
% --------------------------
\begin{mybox}
\textbf{Principle.} Whenever a system is decomposed into subsystems, the deviation from local Fisher additivity defines an intrinsic curvature that encodes correlations, entanglement, or constraints.
\end{mybox}

Let $g$ denote the Fisher information metric on a statistical manifold. For a partition $A|B$,
\[
\mathcal{F} := g^{AB} - (g^A \oplus g^B).
\]

- $\mathcal F = 0$ $\;\Rightarrow$ independent parts, flat geometry.
- $\mathcal F > 0$ $\;\Rightarrow$ non-factorization, excess information, curvature.

\noindent The scalar contraction $\mathcal R_F := \mathrm{tr}(g^{-1}\mathcal F)$ is the \emph{Fisher curvature}.

% --------------------------
\section*{2. Minimal Equations}
% --------------------------
\subsection*{2.1 Additivity baseline}
\[
\mathcal I_{A\cup B} = \mathcal I_A + \mathcal I_B.
\]

\subsection*{2.2 Excess Fisher curvature}
\[
\mathcal I_{A\cup B} = \mathcal I_A + \mathcal I_B + \mathrm{tr}(g^{-1}\mathcal F).
\]

\subsection*{2.3 Classical decomposition}
For score functions $S_A,S_B$ and residual $R$:
\[
\mathcal I_{AB} = \mathcal I_A + \mathcal I_B + 2\,\mathrm{Cov}(S_A,S_B) + \mathrm{Var}(R).
\]

\subsection*{2.4 Quantum scaling}
\begin{align*}
\text{Separable: } & I(\theta) \sim N, \\
\text{Entangled: } & I(\theta) \sim N^2.
\end{align*}

\subsection*{2.5 Holographic duality}
\[
\text{Boundary Fisher metric} = \text{Bulk canonical energy}.
\]

% --------------------------
\section*{3. Operational Heuristics}
% --------------------------
\begin{enumerate}[label=(\alph*)]
    \item Estimate KL or Fisher across partitions: non-additivity $\Rightarrow$ excess curvature.
    \item If boundary relative entropy is preserved: interior measure-conjugacy is possible.
    \item Divergence of $\mathcal R_F$: signals criticality or phase transitions.
\end{enumerate}

% --------------------------
\section*{4. Cross-Domain Examples}
% --------------------------
\subsection*{4.1 Physics}
- Entangled probes yield excess QFI beyond separable limits.
- Black-hole entropy area law: entropy $\sim$ surface, not volume. Boundary Fisher curvature encodes bulk geometry.

\subsection*{4.2 Neuroscience}
- Neural codes exhibit synergy: joint Fisher info exceeds sum of parts.
- Integrated information measures can be reframed as excess Fisher curvature.

\subsection*{4.3 Economics}
- Market co-movements correspond to spikes in Fisher curvature.
- Stress periods reveal super-additivity in covariance structure.

\subsection*{4.4 Climate/Earth systems}
- Teleconnections (e.g. ENSO) manifest as $\mathcal R_F > 0$ across spatial regions.
- Critical transitions marked by curvature divergence.

\subsection*{4.5 Machine Learning}
- Feature interactions: positive $\mathcal F$ indicates synergy not captured by independent predictors.

% --------------------------
\section*{5. Programmatic Steps}
% --------------------------
\begin{enumerate}[label=Step \arabic*.]
    \item \textbf{Calculus of cuts:} develop transformation rules for $\mathcal F$ under coarse-graining and embeddings.
    \item \textbf{Curvature budgets:} partition observed $\mathcal F$ into noise vs hidden geometry vs hidden drivers.
    \item \textbf{Design by curvature:} maximize useful $\mathcal F$ (metrology) or minimize harmful $\mathcal F$ (stability).
    \item \textbf{Correspondence library:} map domain-specific invariants (Bell, JLMS, Ruppeiner) to contractions of $\mathcal F$.
\end{enumerate}

% --------------------------
\section*{6. Outlook}
% --------------------------
The IGC provides a unifying framework: just as calculus gave a universal language for infinitesimal change, IGC proposes a universal language for \emph{non-additivity}. Fisher curvature is its central invariant, linking statistics, quantum mechanics, and spacetime geometry.

\end{document}
